% Junior Level 2, Week 8 -- workbook chapter (read aloud by a grown-up)
\chapterweek{8}{Start, Win and Lose}{A game needs a start, a goal and an ending.}

\section{What you will learn}
\begin{ican}
  \item I can give my game a start and a goal.
  \item I can make my game end: you win, or game over.
  \item I can fix a game that never ends.
\end{ican}

\section{Three parts of a game}
A game has a \textbf{start}, a \textbf{play} part and an \textbf{ending}. The block \enquote{stop all} is the final whistle.

\begin{exercise}[Win and Lose]
\begin{enumerate}[leftmargin=1.4em,itemsep=2pt]
  \item Open your catch game. Make a variable called \enquote{misses}.
  \item Apple: put \enquote{change misses by 1} in its edge check.
  \item Bowl: build the referee script.
  \item Play to win. Then play to lose!
\end{enumerate}
\begin{center}
\begin{scratch}[scale=0.8]
\blockinit{when \greenflag clicked}
\blockvariable{set \selectmenu{score} to \ovalnum{0}}
\blockvariable{set \selectmenu{misses} to \ovalnum{0}}
\blocklook{say \ovalnum{Catch 10 apples!} for \ovalnum{2} seconds}
\blockinfloop{forever}{
  \blockif{if \booloperator{\ovalvariable{misses} = \ovalnum{3}} then}{
    \blocklook{say \ovalnum{Game over!} for \ovalnum{2} seconds}
    \blockstop{stop \selectmenu{all}}
  }
  \blockif{if \booloperator{\ovalvariable{score} = \ovalnum{10}} then}{
    \blocklook{say \ovalnum{You win!} for \ovalnum{2} seconds}
    \blockstop{stop \selectmenu{all}}
  }
}
\end{scratch}
\end{center}
\end{exercise}

\needspace{26\baselineskip}
\begin{exercise}[Never or Right Away?]
Leo's game ends right at the start. Zoe catches 10 apples, but her game never ends. Circle the bug in each.

\begin{center}
\begin{tabular}{@{}c@{\hspace{1cm}}c@{}}
\textbf{Leo} & \textbf{Zoe}\\[2pt]
\begin{scratch}[scale=0.7]
\blockinit{when \greenflag clicked}
\blockvariable{set \selectmenu{misses} to \ovalnum{0}}
\blockinfloop{forever}{
  \blockif{if \booloperator{\ovalvariable{misses} = \ovalnum{0}} then}{
    \blocklook{say \ovalnum{Game over!} for \ovalnum{2} seconds}
    \blockstop{stop \selectmenu{all}}
  }
}
\end{scratch}
&
\begin{scratch}[scale=0.7]
\blockinit{when \greenflag clicked}
\blockvariable{set \selectmenu{score} to \ovalnum{0}}
\blockif{if \booloperator{\ovalvariable{score} = \ovalnum{10}} then}{
  \blocklook{say \ovalnum{You win!} for \ovalnum{2} seconds}
  \blockstop{stop \selectmenu{all}}
}
\blockinfloop{forever}{
  \blockif{if \booloperator{\ovalvariable{misses} = \ovalnum{3}} then}{
    \blocklook{say \ovalnum{Game over!} for \ovalnum{2} seconds}
    \blockstop{stop \selectmenu{all}}
  }
}
\end{scratch}
\end{tabular}
\end{center}
Try each bug in a copy of your game. Then fix it.
\end{exercise}

\begin{exercise}[My Game Rules]
Next week you make your own arcade game! Circle one: \quad catch game \quad maze game
\begin{itemize}[leftmargin=1.4em,itemsep=6pt]
  \item Start: the player learns \rule{6cm}{0.4pt}
  \item I win when \rule{7cm}{0.4pt}
  \item I lose when \rule{7cm}{0.4pt}
\end{itemize}
\end{exercise}

\begin{vocab}
  \vocabword{goal}{What you try to do to win, like catch 10 apples.}
  \vocabword{game over}{The ending when you lose. The game stops.}
\end{vocab}

\begin{reflection}
Color the parts your game has: \quad \fbox{\strut\ start\ } \fbox{\strut\ play\ } \fbox{\strut\ ending\ }

How do you feel today? Circle one.

\begin{center}
\Huge\faSmile[regular]\qquad\faMeh[regular]\qquad\faFrown[regular]
\end{center}
\end{reflection}
